\scanpage{27}
i.e.
\[\boxed{\Phi'=\Phi+(pL-F_S^*(L))}\]
On voit de même, si $\omega=\omega_{M,\varphi}\in\Gamma(\widehat\Omega^1_{S/W}\otimes(T\otimes_{\mathbb Z_p}\mathcal O_S))$ donne la forme connexion $\nabla\in\Gamma(\widehat\Omega^1_{S/W}\otimes M)$ via $\varphi$,
\[
\nabla_{M,\varphi}=\begin{pmatrix}0&\omega_{M,\varphi}\\0&0\end{pmatrix}
\]
\margin{$\boxed{\omega=\omega_{M,\varphi}\in\Gamma(\widehat\Omega^1_{S/W}\otimes(T\otimes_{\mathbb Z_p}\mathcal O_S))}$}
qu'on a, si $\omega'=\omega_{M,\varphi'}$, \unsure{qu'on a}
\[\boxed{\omega'_{\deleted{\unsure{M,\varphi}}}=\omega-dL}\]

Écrivons les conditions \deleted{\unsure{qui}} sur \deleted{\unsure{un tel}} $\Phi=\Phi_{M,\varphi}$ et $\omega=\omega_{M,\varphi}$ qui assurent \unsure{qu'on a} bien un objet de $C$. Ce sont
\[
\left\{\begin{array}{l}
\boxed{d\omega=0}\\[1em]
\boxed{F_S^*(\omega)-p\omega=d\Phi}\\[1em]
\boxed{\Phi_0=0}
\end{array}\right.
\]
\margin{exprime que la connexion $\nabla$ définie par $\unsure{\varphi|\omega}$ est à courbure nulle\par
exprime que $F_M$ est compatible à la connexion\par
exprimant que $(F_M)_0$ respecte \deleted{\unsure{le}} $\mathrm{Fil}_0$.}
\mspage{26}
